\section{Preliminaries}\label{section2}
In this section we recall some basic concepts and notations from \cite{jps} and then we shall prove some new preliminary results, which are needed later on.
	
\subsection{Connected (co)rings} \label{corings} Let $R$ be a semisimple ring that we f\/ix throughout the article. Since we will always work with algebras and coalgebras in the tensor category of $R$-bimodules, an unadorned tensor product $\ot$ will mean $\ot_R$. Let $V$ be an $R$-bimodule. The notation $V^{(n)}$ will be used for the $n$-th tensor power $V \ot \cdots \ot V$. Conventionally, $V^{(0)} = R$. Similarly, for any bimodule morphism $f\colon M \to N$, the tensor product $f \ot \cdots \ot f$ with $n$ factors will be denoted by $f^{(n)}$. For a set $X$, the identity morphism will be denoted by $\operatorname{Id}_X$ or simpler, $\operatorname{I}_X$ or even $X$ when there is no risk of confusion.

An $\bb{N}$-graded algebra $A = \op_{n \in \bb{N}} A^n$ is called an \emph{$R$-ring} if it is an algebra in the tensor category of $R$-bimodules and it is connected if $A^0=R$. Dually, a connected $R$-coring is a graded coalgebra $C = \op_{n \in \bb{N}}$ in the category of $R$-bimodules such that $C_0=R$. The augmentation ideal of $A$ will be denoted by $A_+ = \op_{n \in \bb{N}^\ast} A^n$ and similarly for~$C$. Note that the comultiplication map $\Delta$ induces a canonical coassociative map $\Delta_+ \colon C_+ \to C_+ \ot C_+$, where $C_+$ can also be viewed as~$C/R$.

For a graded connected $R$-ring $A$, the multiplication $\mu$ is def\/ined by some maps $\mu^{p,q} \colon$ \mbox{$A^p \ot A^q \to A^{p+q}$}. When all of these maps are surjective, for all $p,q \geq 0$, we say that $A$ is \emph{strongly graded}. Equivalently, $A$ is strongly graded if and only if the iterated multiplication $\mu_n \colon (A^1)^{(n)} \to A^n$ is surjective for all $n \geq 2$.

Similarly, for a graded connected $R$-coring $C$, the comultiplication is def\/ined by some $R$-bimodule maps $\Delta_{p,q} \colon C_{p+q} \to C_p \ot C_q$ and the coring $C$ is called \emph{strongly graded} if and only if all of them are injective.

Let $\Delta^1 = \operatorname{Id}_{C_1}$ and for $n \geq 2$ def\/ine a map $\Delta^n\colon C_n \to C_1^{(n)}$ by the recurrence relation
\begin{gather*}
	\Delta^n = \big(\operatorname{Id}_{C_1} \ot \Delta^{n-1}\big) \circ \Delta_{1,n-1}.
\end{gather*}

It is immediate from the def\/inition and the coassociativity property that the following equality holds true
\begin{gather*}
	\Delta^{p+q} = (\Delta^p \ot \Delta^q) \circ \Delta_{p,q}.
\end{gather*}
Another remark is that the strong grading on $C$ can be seen to be equivalent to the injectivity of all of the maps $\Delta^n$, $\forall\, n$ and further that $\Delta_{1,n}$ is injective for all $n$ if and only if $\Delta_{n,1}$ is injective for all~$n$.

When $C$ is a connected $R$-coring, the unit of $R$ is a group-like element, that is $\Delta(1) = 1 \ot 1$, cf.~\cite{jps}. Therefore, we can speak of primitive elements of~$C$, namely those $c \in C$ for which $\Delta(c) = 1 \ot c + c \ot 1$. The set of all primitive elements in $C$ contains $C_1$ and will be denoted by~$PC$. In general, the inclusion of $C_1$ in $PC$ is strict, but as per \cite[Lemma~1.4]{mst}, $PC = C_1$ if and only if~$C$ is strongly graded.

Indecomposable elements of an augmented $R$-ring correspond by duality to primitive elements. They were introduced by May in~\cite{may} and we will see how they relate to the strong grading on an $R$-ring. In the graded case, the $R$-bimodule $QA$ of indecomposable elements is def\/ined by the exact sequence
\begin{gather*} A_+ \ot A_+ \xrar{\mu} A_+ \to QA \to 0.\end{gather*}
There is a canonical morphism $A^1 \to QA$, which maps $a \in A^1$ to its class $a + A_+^2 \in QA$. This map has a left inverse, induced by the projection $A_+ \to A^1$.

There is a canonical morphism $A^1 \to QA$ which maps an element in $A^1$ to its class modulo the square of the augmentation ideal $A_+$ in~$QA$. This map has a left inverse, induced by the projection $A_+ \to A^1$.

We can prove a f\/irst result, which is a dual version of \cite[Lemma~1.4]{mst} and which shows the role of $QA$, the bimodule of indecomposable elements in an $R$-ring. Moreover, the following lemma will be used for obtaining our main characterization theorem for Koszul corings (i.e., Theorem~\ref{thm:kcoring}).

\begin{Lemma} \label{lema:ring}
Let $A$ be a graded and connected $R$-ring.
\begin{enumerate}\itemsep=0pt
\item[$1.$] $A$ is strongly graded if and only if the canonical map $QA\to A^1$ is injective, if and only if the canonical map $A^1\to QA$ is surjective.

\item[$2.$] If $A$ is strongly graded, $B$ is a connected graded $R$-ring and $g\colon B \rar A$ is a morphism of graded $R$-rings such that its components $g^0$ and $g^1$ are surjective, then $g$ is surjective.

\item[$3.$] Let $A=\oplus_{n,m\in\bb{N}}A^{n,m}$ be bigraded. If $\operatorname{gr}A$ is strongly graded and $A^{n,m}=0$ for $n=0,1$ and all $m \neq n$, then $A^{n,m}=0$ for all $n\geq 2$, $m\neq n$.
\end{enumerate}
\end{Lemma}

\begin{proof}
Since $QA\to A^1$ is a left inverse of $A^1\to QA$ it follows that the former map is injective if and only if the latter is surjective. On the other hand, $A$ is strongly graded if and only if $A_n\subseteq A_+^2$, for all $n\geq 2$. Consequently, the f\/irst statement holds, as the above inclusions are true if and only if the map $A^1\to QA$ is surjective.

The proof of the second statement follows as in the case of corings \cite[Lemma~1.4]{mst}, using the diagram below
\begin{gather*}
 \xymatrixcolsep{2pc}\xymatrix{\ar[r]^-{\mu^{n,1}_{B}} B_{n} \otimes B_{1} \ar[d]_-{g_n \otimes g_1} &B_{n+1} \ar[d]^-{g_{n+1}} \\
A_{n} \otimes A_{1} \ar[r]^-{\mu^{n,1}_{A}} & A_{n+1}. }
\end{gather*}
For proving the last part of the lemma we def\/ine $\operatorname{Diag}(A)$ to be graded subring of $\operatorname{gr}(A)$ given by $\operatorname{Diag}(A)=\oplus_{n\in\N} A^{n,n}$. Let $\iota\colon \operatorname{Diag}(A)\to \operatorname{gr}(A)$ denote the inclusion map. Obviously, by the standing assumptions, $\iota^0$ and $\iota^1$ are surjective maps. Hence, by the second part of the lemma, $\iota$~is surjective. Thus $A^{n,m}$ must be zero for all $n$ and $m$, with $n\neq m$.
\end{proof}

\subsection[The $R$-ring $\langle V,W\rangle$ and the $R$-coring $\{V,W\}$]{The $\boldsymbol{R}$-ring $\boldsymbol{\langle V,W\rangle}$ and the $\boldsymbol{R}$-coring $\boldsymbol{\{V,W\}}$}\label{shriek}

For an $R$-bimodule $V$ and a sub-bimodule $W\subseteq V\ot V$, we def\/ine the $R$-ring $\langle V, W \rangle$ to be the quotient $T^a_R (V)/\langle W \rangle$ of the tensor algebra of the $R$-bimodule $V$ by the two-sided ideal generated by $W$. Note that $\langle W\rangle=\sum\limits_{n\in\mathbb{N}} \langle W\rangle^n$, where
 \begin{gather*}
 \langle W\rangle^n=\sum_{i=1}^{n-1}V^{(i-1)}\ot W\ot V^{ (n-i-1)}. \end{gather*}
For $V$ and $W$ as above, one also constructs a graded $R$-coring $\br{V,W}$ by taking $\br{V,W}_0=R$ and $\br{V,W}_1=V$. For all $n \geq 2$ def\/ine
\begin{gather*}
\br{V,W}_n=\bigcap_{p=1}^{n-1} V^{(p-1)} \otimes W \otimes V^{ (n-p-1)}.
\end{gather*}
As it is shown in~\cite{jps}, the direct sum $\br{V,W}=\oplus_{n \in \mathbb{N}} \br{V,W}_n$ is a graded subcoring of $T^c_R(V)$, the tensor coalgebra of the $R$-bimodule~$V$.

If $A$ is a connected graded $R$-ring and $C$ is a connected graded $R$-coring then the graded coring $\{A^1,\Ker\mu^{1,1}\}$ and the graded ring $\gen{C_1,\im\Delta_{1,1}}$ will be denoted by $A^!$ and $C^!$, respectively and called the \emph{shriek coring} and the \emph{shriek ring}, respectively. The homogeneous component of degree~$n$ of~$A^!$ will be denoted by $A^!_n$. To simplify the notation, for the ring $C^!$ we shall write $C^!_n$ instead of~$(C^!)^n$.

\begin{Remark} In the literature, the shriek construction (also known as \emph{the quadratic dual} \cite[Section~2.8]{bgs}, for example) is def\/ined and used slightly dif\/ferent that our setting. However, we note that throughout this paper, the only meaning of the shriek structures is that introduced above. Note that, unlike the classical case of~\cite{bgs}, for example, here the shriek construction changes an $R$-ring into an $R$-coring and viceversa. \end{Remark}

\subsection{Bigraded corings}
Some basic facts regarding bigraded corings are due, since many of the structures which we will use are of this kind.

An $R$-coring $C$ is called \emph{bigraded} if it has a decomposition as a direct sum of $R$-bimodules $C=\oplus_{n,m \in \mathbb{N}} C_{n,m}$ such that its comultiplication induces a collection of maps
\begin{gather*} C_{n+n',m+m'} \xrightarrow{\Delta_{n,m}^{n',m'}} C_{n,m} \otimes C_{n',m'}.\end{gather*}
For this case, coassociativity means that the diagram below is commutative, for all positive integers $m$, $n$, $p$, $m'$, $n'$ and $p'$,
\begin{gather*}
 \xymatrixcolsep{5pc}\xymatrix{ C_{n+m+p,n'+m'+p'} \ar[d]_-{\Delta_{n,n'}^{m+p,m'+p'}} \ar[r]^-{\Delta_{n+m,n'+m'}^{p,p'}} & C_{n+m,n'+m'} \otimes C_{p,p'} \ar[d]^-{\Delta_{n,n'}^{m,m'} \otimes \Id_{C_{p,p'}}} \\
C_{n,n'} \otimes C_{m+p,m'+p'} \ar[r]_-{\Id_{C_{n,n'}} \otimes \Delta_{m,m'}^{p,p'}} & C_{n,n'} \otimes C_{m,m'} \otimes C_{p,p'}.}
\end{gather*}
By def\/inition, we also impose to the counit to vanish on $C_{n,m}$, provided that either $n>0$ or $m>0$.

Starting with a bigraded coring, one can associate to it a graded coring $\operatorname{gr}(C)$, whose homogeneous component of degree $n$ is $\operatorname{gr}_n(C)=\oplus_m C_{n,m}$. Therefore, def\/ine $\operatorname{gr}(C):=\oplus_n \operatorname{gr}_n(C)$. We will be interested only in \textit{connected bigraded corings}, namely those bigraded corings for which $C_{0,0}=R$ and $C_{0,m}=0$, for $m>0$. Note that $C$ to be connected implies $\operatorname{gr}(C)$ is connected as well.

Let $C$ be a (connected) bigraded coring. Def\/ine
\begin{gather*} C'_{n,m}=\begin{cases} C_{n,m}, & n=m, \\ 0, & n\neq m.\end{cases}\end{gather*}
Then ${C'}:=\oplus_{n,m} {C}_{n,m}'$ becomes a (connected) bigraded coring. We denote the graded coring $\operatorname{gr}({C'})$ associated to it by $\operatorname{Diag}(C)$.

Keeping the notations and the context above, there exist canonical $R$-bimodule morphisms $\pi_{n,m} \colon C_{n,m} \to C'_{n,m}$ which are identities on $C_{n,n}$ and zero maps in rest. Their collection, $\pi = \{\pi_{n,m}\}_{n,m}$, def\/ines a morphism in the category of graded corings. Since $\operatorname{Diag}(C) = \operatorname{gr}(C')$, it follows that the map $\pi$ induces a morphism at the level of graded corings, namely $\operatorname{gr}(\pi) \colon \operatorname{gr}(C) \to \operatorname{Diag}(C)$. Surely the homogeneous component of degree $n$ of $\operatorname{Ker}\pi$ is $\op_{n \neq m} C_{n,m}$.

Note that one can def\/ine by duality the corresponding notions of bigraded $R$-rings. We omit the details here, since this case is better known than the one for corings. See, for example, \cite{pi} for the dif\/ferential graded setting.

\subsection{The normalized (co)chain complex} \label{sec-t(a)}
We shall compute $T_n(A):=\Tor_n^A(R,R)$ as the $n$th homology group of the normalized bar complex $(\Omega_\bullet(A), d_\bullet)$, where $\Omega_n(A)=\Ap{n}$. The morphisms $d_n\colon \Omega_n(A)\to\Omega_{n-1}(A)$ are def\/ined by $d_1=0$ and, for $n>1$,
\begin{gather*}
d_n(a_1\ot\cdots \ot a_n)=\sum_{i=1}^{n-1}(-1)^{i-1}a_1\ot\cdots \ot a_ia_{i+1}\ot\cdots\ot a_n.
\end{gather*}
Since the normalized bar complex has a canonical structure of DG-coalgebra in the category of $R$-bimodules with respect to the comultiplication $\Delta_{p,q}(a_1\ot\cdots \ot a_{p+q})=(a_1\ot\cdots \ot a_{p})\ot(a_{p+1}\ot\cdots \ot a_{p+q})$ it follows that $T(A)=\oplus_{n \in \mathbb{N}} T_n(A)$ has a canonical structure of connected $R$-coring.
	
One can express the complex $\Omega_\bullet(A)$ as a direct sum $\Omega_\bullet(A)=\oplus_{m \geq 0} \Omega_\bullet(A,m)$ of subcomplexes. As in \cite[Section~1.6]{mst}, we introduce the following terminology and notation. An $n$-tuple $\boldsymbol{m}=(m_1,\dots,m_n)$ is called \textit{a positive $n$-partition of $m$} if and only if $\sum\limits_{i=1}^n m_i=m$ and all $m_i$ are positive. The set of all positive $n$-partitions of $m$ will be denoted by $\mathscr{P}_n(m)$. Furthermore, if $\boldsymbol{m}=(m_1,\dots,m_n)$ is a positive $n$-partition and $A$ is a connected $R$-ring then for the tensor product $A^{m_1}\ot\cdots\ot A^{m_n}$ we shall use the notation $A^{\boldsymbol{m}}$. For a positive $n$-partition $\bs{m}=(m_1,\dots,m_n)$ of~$m$, the multiplication $\m$ of $A$ induces bimodule maps $\mu^{\bs{m}}\colon A^{\bs{m}}\to A^m$ and $\mu_{\bs{m}}\colon (A^1)^{(m_1)}\ot\cdots \ot(A^1)^{(m_n)}\to A^{\bs{m}}$. Note that, by def\/inition, $\mu_{\bs{m}}=\m(m_1)\ot\cdots\ot\m(m_n)$.

Hence, with this notation, $\Omega_\bullet(A,m)$ is the following subcomplex of $\Omega_\bullet(A)$:
\begin{gather*}
0 \xleftarrow{d_0^m} 0 \xleftarrow{d_1^m} \bigoplus \limits_{\boldsymbol{m}_1\in\mathscr{P}_1(m)} A^{\boldsymbol{m}_1} \xleftarrow{d_2^m} \bigoplus \limits_{\boldsymbol{m}_2\in\mathscr{P}_2(m)} A^{\boldsymbol{m}_2}\xleftarrow{d_3^m} \cdots \xleftarrow{d^m_n}  \bigoplus \limits_{\boldsymbol{m}_n\in\mathscr{P}_{n}(m)} A^{\boldsymbol{m}_n}\xleftarrow{\ \ \ }\cdots .
\end{gather*}
Of course, there is only one $1$-partition of $m$, namely $\boldsymbol{m}_1=(m)$. Thus the f\/irst direct sum in~$\Omega_\bullet(A,m)$ coincides with $A^m$. The homology in degree $n$ of $\Omega_\bullet(A,m)$ will be denoted either by~$T_{n,m}(A)$ or~$\operatorname{Tor}_{n,m}^A(R,R)$. Clearly, we have $T(A)=\oplus_{m\geq 0} T_{n,m}(A)$ and this decomposition is compatible with the coring structure, in the sense that $T(A)$ is a bigraded coring with~$T_{n,m}(A)$ as $(n,m)$-homogeneous component.
	
Dually, for a connected $R$-coring $C$, the normalized bar cochain complex $(\Ou{\bullet}{C}, d^\b)$ is def\/ined by $\Ou{n}{C}=\Cp{n}$ and $d^0=0$, while
\begin{gather*}
d^n=\sum_{i=1}^{n}(-1)^{i-1}\I_{C_+^{(i-1)}}\ot\;\Delta_+\ot\I_{C_+^{(n-i)}}.
\end{gather*}
Here $C_+:=C/C_0$ and $\Delta_+\colon C_+\to C_+\ot C_+$ is the map induced by the comultiplication of~$C$. It is well-known that $E^n(C)=\operatorname{Ext}^n_C(R,R)$ is the $n$th cohomology group of $\Omega^\bullet(C)$ and that $E(C)=\oplus_{n \in \mathbb{N}} \ E^n(C)$ is a connected $R$-ring with respect to the multiplication induced by the DG-algebra structure (in the category of $R$-bimodules) on $\Ou{\b}{C}$ which is def\/ined by concatenation of tensor monomials (see, e.g., \cite[Section~1.1]{PP}).

The complex $\Omega^\b(C)$ is a direct sum of subcomplexes $\Omega^\b(C,m)$, that are def\/ined as follows. For a positive $n$-partition of $m$ let $C_{\boldsymbol{m}}:=C_{m_1}\ot \cdots\ot C_{m_n}$. Hence $\Omega^\b(C,m)$ is the subcomplex
\begin{gather*}
\ 0 \xrightarrow{\ \ \ } 0 \xrightarrow{d_m^0} \bigoplus \limits_{\boldsymbol{m}_1\in\mathscr{P}_1(m)} C_{\boldsymbol{m}_1} \xrightarrow{d_m^1} \bigoplus \limits_{\boldsymbol{m}_2\in\mathscr{P}_2(m)} C_{\boldsymbol{m}_2} \xrightarrow{d_m^2} \cdots \xrightarrow{d^{n-1}_m}  \bigoplus \limits_{\boldsymbol{m}_n\in\mathscr{P}_{n}(m)} C_{\boldsymbol{m}_n} \xrightarrow{d_m^n}\cdots .
\end{gather*}
Of course, $\Omega^1(C,m)=C_m$. The homology in degree $n$ of $\Omega^\bullet(C,m)$ will be denoted either by $E^{n,m}(C)$ or $\operatorname{Ext}^{n,m}_C(R,R)$. Clearly, we have $E(C)=\oplus_{m\geq 0}E^{n,m}(C)$ and this decomposition is compatible with the ring structure in the sense that $E(C)$ is a bigraded ring with $E^{n,m}(C)$ as $(n,m)$-homogeneous component. For more details the reader is referred to \cite[Section~1.15]{jps}.

\subsection{Almost-Koszul pairs} \label{sec:almost}
As mentioned in the introduction, an almost Koszul pair $(A,C)$ consists of a connected and graded $R$-ring $A$ and a connected graded $R$-coring $C$, together with an isomorphism of $R$-bimodules $\theta_{A,C} \colon C_1 \rightarrow A^1$ which satisf\/ies the relation \eqref{ec:theta}. Using the graded version of Sweedler's notation for the comultiplication of $C$, the above equation is equivalent to
\begin{gather} \label{ec:qksz}
	\sum \theta_{C,A}(c_{1,1})\theta_{C,A}(c_{2,1})=0,
\end{gather}
where $c$ is an arbitrary element of $C_2$.

If $(A,C)$ and $(B,D)$ are almost-Koszul pairs, then a morphism of almost-Koszul pairs is a~couple $(\phi,\psi)$, where $\phi\colon A \rar B$ is a morphism of graded $R$-rings and $\psi\colon C \rar D$ is a morphism of graded $R$-corings such that they commute with the isomorphisms $\theta_{A,C}$ and $\theta_{B,D}$. Henceforth, the diagram below is commutative,
\begin{gather*}
 \xymatrix{ A^1 \ar[r]^-{\phi_1} \ar[d]_-{\theta_{A,C}} & B^{1} \ar[d]^-{\theta_{B,D}} \\
C_1 \ar[r]^-{\psi_1} & D_1.}
\end{gather*}
A f\/irst example of an almost-Koszul pair is given in \cite[Proposition 1.8]{jps}. For any connected and strongly graded $R$-ring $A$, the pair $\big(A,T(A)\big)$ is almost-Koszul. Here, the coring structure of $T(A)$ is def\/ined as in Section~\ref{sec-t(a)}, and the map $\theta_{A,T(A)}\colon T_1(A)\to A^1$ is induced by the projection $A_+\to A^1$. Note that $T_1(A)=A_+/A_+^2\cong A^1$, as $A$ is strongly graded, so $\theta_{A,T(A)}$ is an isomorphism.

The couple $(A,A^!)$ is another example of an almost-Koszul pair. Recall that $A^!_1=A^1$, so we can take $\theta_{A,A^!}:=\Id_{A^1}$. The condition~\eqref{ec:theta} is trivial in this particular case as, by construction, $A^!_2=\Ker\mu^{1,1}$.

Dually, if $C$ is a connected and strongly graded $R$-coring, the pair $(E(C),C)$ is almost-Koszul, cf.~\cite[Proposition~1.18]{jps}. In this example the $R$-ring structure of $E(C)$ is def\/ined in Section~\ref{sec-t(a)}. Since $E^1(C)=\Ker\Delta_+$, the map $\theta_{E(C),C}\colon C_1\to E^1(C)$, given by $\theta_{E(C),C}(c):=c+C_0\in C/C_0$, is well def\/ined and it is an isomorphism of $R$-bimodules.

The couple $(C^!,C)$ can be seen as an almost-Koszul pair with respect to the map $\theta_{C^!,C}=\Id_{ C_1}$. The relation \eqref{ec:theta} is verif\/ied since $C^!_2:=(C_1\ot C_1)/\im\Delta_{1,1}$ and the component $C_1^!\ot C_1^!\to C^!_2$ of the multiplication on $C^!$ coincides with the canonical projection from $C_1\ot C_1$ onto $C^!_2$.

\subsection{Koszul pairs} \label{fa:Koszul pairs}
Following \cite{jps} we shall brief\/ly recall the def\/inition of Koszul pairs, which are our main tool to investigate Koszul $R$-(co)rings. For any almost-Koszul pair $(A,C)$ and $n\geq 0$ we def\/ine a~complex of graded right $C$-comodules by
\begin{gather*} \operatorname{K}_r^{-1}(A,C)=R \qquad \text{and} \qquad \operatorname{K}^n_r(A,C)=A^n \otimes C, \qquad \forall\, n \geq 0.\end{gather*}
The dif\/ferential map $d_{r}^{n}\colon A^{n}\otimes C\rightarrow A^{n+1}\otimes C$ is zero on $A^{n}\otimes C_{0}$ and, for $p>0$ and $a\otimes c\in A^{n}\otimes C_{p},$
\begin{gather*}
 d_{r}^{n}(a\otimes c)=\sum\limits a\theta _{C,A}(c_{1,{1}})\otimes c_{2,{p-1}}.
\end{gather*}
The dif\/ferential $d^n_r$ maps $A^{n}\ot C_{p}$ to $A^{n+1}\ot C_{p-1}$. Thus $\K^\bullet_r(A,C,m)$ is a subcomplex of $\K^\bullet_r(A,C)$, for all $m\in\N$, where $\K^n_r(A,C,m)= A^{n}\ot C_{m-n}$. Note that, by convention, $C_p=0$ for $p<0$, so $\K^n_r(A,C,m)$ is trivial if either $n<-1$ or $n>m$.

Using the fact that $(A^{\rm op},C^{\rm op})$ is an almost-Koszul pair over $R^{\rm op}$, cf.\ \cite[Remark~1.4]{jps}), a~complex of left $C$-comodules is obtained by setting $\operatorname{K}_l^\bullet(A,C)=\operatorname{K}_r^\bullet(A^{\rm op},C^{\rm op})$.

By combining the complexes $\operatorname{K}_l^\bullet(A,C)$ and $\operatorname{K}_r^\bullet(A,C)$, in \cite{jps} one constructs another cochain complex $\operatorname{K}^\bullet(A,C)$ in the category of $C$-bicomodules. Since we do not use it in this paper, we omit its def\/inition.

By duality, to an almost-Koszul pair correspond also three chain complexes $\operatorname{K}^l_\bullet\!(A{,}C)$, $\operatorname{K}^r_\bullet\!(A{,}C)$ and $\operatorname{K}_\bullet(A,C)$. For instance, $\operatorname{K}^l_\bullet(A,C)$ is the complex of left $A$-modules,
\begin{gather*} \operatorname{K}_{-1}^l(A,C)=R \qquad \text{and} \qquad \operatorname{K}^l_n(A,C)=A \ot C_n,\end{gather*}
whose dif\/ferential $d_{0}^l$ is given by the left action of $A$ on $R=C_0$. For $n > 0$, the map $d_n^l$ acts as
\begin{gather*}
	d_{n}^{l}(a\otimes c)=\sum\limits a\theta _{C,A}(c_{1,{1}})\otimes c_{2,{n-1}}.
\end{gather*}
The complex $\operatorname{K}_\bullet^l(A,C)$ also decomposes as a direct sum $\oplus_{m\geq 0}\operatorname{K}^l_\bullet(A,C,m)$ of subcomplexes, where ${K}^l_n(A,C,m)=A^{m-n}\ot C_n$. Note that, ${K}^l_n(A,C,m)=0$, for $n>m$.

We conclude this subsection by recalling that the exactness of any of these six complexes implies the exactness of all the others, cf.\ \cite[Theorem~2.3]{jps}. Whenever this is the case, $(A,C)$ is called \emph{Koszul}. Thus, for a Koszul pair $(A,C)$, the complexe $\operatorname{K}_\bullet^l(A,C)$ is a resolution of $R$ by graded projective left $A$-modules. Similarly, for such a pair, $\operatorname{K}^\bullet_r(A,C)$ is a resolution of $R$ by injective graded right $C$-comodules.

The subcomplex $\K_\bullet^l(A,C,0)$ is trivial in degree $n$, for all $n\neq -1,0$ and $d_{-1}^l\colon R\to A^0\ot C_0$, coincides with the canonical isomorphism $R\cong R\ot_R R$. Thus, $\K_n^l(A,C,0)$ is always exact. We deduce that the pair $(A,C)$ is Koszul if and only if the complexes $\K_\bullet^l(A,C,m)$ are exact for all $m>0$. By duality, $(A,C)$ is Koszul if and only if $\K^\bullet_r(A,C,m)$ is Koszul for all $m>0$.

\subsection{Examples of morphisms of almost-Koszul pairs}\label{fa:exemples_morphisms}
We start by def\/ining a morphism $(\Id_A,\phi^A)$ from $(A,A^!)$ to $\big(A,T(A)\big)$. The graded coring morphism $\phi^A$ is obtained applying \cite[Proposition 1.24]{jps} as follows. We keep the notation from the above mentioned result. Deleting the component of degree $-1$ and then applying the functor $R\ot_A (-)$ to $\phi _{\bullet}\colon \K_{\bullet}^{l}(A,A^!)\to \beta_{\bullet}^{l}(A)$ we get a morphism of complexes from $R\ot_A \K_{\bullet}^{l}(A,A^!)$ to $R\ot_A \beta_{\bullet}^{l}(A)$. Note that the last two complexes are concentrated in non-negative degrees. By \cite[Proposition 1.23]{jps} the former complex is isomorphic to $(A^!,0)$, the complex with tri\-vial dif\/ferential maps and whose module of $n$-chains coincides with~$A^!_n$. On the other hand, by def\/inition, the latter complex equals $\Od{\bullet}{A}$. Since $\Id_R\ot\phi_\bullet$ is compatible with the coring structure of~$A^!$ and~$\Ou{\bullet}{A}$, by applying the homology functor we get the desired coring map $\phi^A\colon A^!\to T(A)$. Let us remark that
\begin{gather*} A^!_n:=\bigcap_{i=1}^{n-1}\big({A^1}\big)^{(i-1)}\ot\Ker\mu^{1,1}\ot \big({A^1}\big)^{(n-i-1)}\subseteq A_+^{(n)},
\end{gather*}
so any $x$ in $A^!_n$ is an $n$-cycle in $\Od{\bullet}{A}$ and $\phi^A(x)$ as the homology class of $x$. It is easy to check now that $(\Id_A,\phi^A)$ is a morphism of almost-Koszul pairs.

In a similar way one def\/ines a canonical morphism of almost Koszul-pairs $(\phi_C,\Id_C)$ from $(E(C),C)$ to $(C^!,C)$. Deleting the component of degree $-1$ of $\phi ^{\bullet}\colon \beta^\bullet_{r}(C)\to \K^\bullet_{r}(C^!,C)$ and then applying the functor $\Hom^C(R, -)$ to the resulting map of complexes we get a morphism $\ol{\phi}^{\bullet}$ from~$\Ou{\bullet}{C}$ to the cochain complex with trivial dif\/ferential maps $(C^!,0)$, cf.\ \cite[Propositions~1.23 and~1.24]{jps}. Henceforth, we can def\/ine the coring map $\phi_C\colon E(C)\to C^!$ by $\phi_C=H^\bullet(\ol{\phi}^{\bullet})$. Let us notice that, for any $x_1,\dots,x_n\in C_+$, we have
\begin{gather*}
 \ol{\phi}^{n}(x_1\ot\cdots \ot x_n)=a_1a_2\cdots a_n\in C^!_n,
\end{gather*}
where $a_i= (\theta_{C^!,C}\circ\ol{\pi}_1)(x_i)\in C^!_1$ and $\ol{\pi}_1\colon C_+\to C_1$ denotes the map induced by the projection $C\to C_1$. Thus, $\phi_C$ maps the cohomology class of an $n$-cocyle $\omega$ to $ \ol{\phi}^{n}(\omega)$.
